% ECONOMICS_PROBLEM: {"id":"C79-13","title":"Determinacy of Dadaist Delta games","jel_code":"C79","source_id":"OP-0219"}
% FIRST_POSTED: 2026-10-09
% ATTEMPT_STATUS: unresolved
% ENTRY_KIND: research_attempt
\documentclass[11pt]{article}
\usepackage{amsmath,amssymb,amsthm,geometry,hyperref}
\geometry{margin=1in}
\newtheorem{theorem}{Theorem}
\newtheorem{proposition}[theorem]{Proposition}
\newtheorem{lemma}[theorem]{Lemma}
\newtheorem{corollary}[theorem]{Corollary}
\theoremstyle{definition}
\newtheorem{example}[theorem]{Example}
\newtheorem{remark}[theorem]{Remark}
\newcommand{\R}{\mathbb R}
\newcommand{\N}{\mathbb N}
\DeclareMathOperator{\NE}{NE}
\DeclareMathOperator{\cav}{cav}
\DeclareMathOperator{\supp}{supp}
\title{Determinacy of Dadaist Delta games\\ GPT 6 Astra Ultra}
\author{Research attempt}
\date{9 October 2026}
\begin{document}
\maketitle

\section*{Target and approach}
For a set-indexed family of well-founded combinatorial game forms, a $\Delta$ move changes one coordinate, limits take coordinatewise eventual values, and the designated First player moves first again after each limit. The open question asks determinacy under these rules. We retain the forms and fixed coordinate identities, not only their Conway values. The definition and open question were checked in Lipparini's Definitions 3.1 and 3.3 and Problem 3.4 \cite{lipparini}.

This attempt establishes a general size bound on individual runs, gives a complete strategy proof for paired conjugate forms, and proves that even one fixed countable example has unbounded countable run lengths. The last result identifies precisely why a bound on each individual run does not supply the missing backward induction. None of these is a determinacy theorem for arbitrary unpaired sums; the pairing argument is classical in spirit and no priority is claimed.

\section*{Set-sized termination, without a countdown rank}
For each coordinate $i$, its option relation is well-founded: there is no infinite sequence of successive options starting from $G_i$. A run is indexed by an ordinal, with legal successor moves and the stipulated limits.

\begin{lemma}[Finite coordinate activity]
In every run, each fixed coordinate is moved only finitely many times. At every limit stage each coordinate is consequently eventually constant.
\end{lemma}
\begin{proof}
If coordinate $i$ moved infinitely many times, its move times, being a subset of an ordinal, would have a first time, a second time, and so on. Taking the resulting forms at those first countably many moves would produce an infinite successive-option sequence in $G_i$. This contradicts well-foundedness. At a limit stage there are therefore only finitely many changes in that coordinate before the limit. A finite set of times below a limit ordinal is bounded below it, so after its last change the coordinate is constant.
\end{proof}

\begin{proposition}[A cardinal upper bound]
Work in ZFC and put $\kappa=\max\{|I|,\aleph_0\}$. Every legal run has length less than $\kappa^+$, where $\kappa^+$ is the successor initial ordinal of $\kappa$. Every legal alternating history has a maximal legal extension. For countable $I$, each completed play has countable length, and histories and strategies can be treated as sets.
\end{proposition}
\begin{proof}
Label each move by the coordinate $i$ it changes and the positive integer counting how many times that coordinate has previously been changed, including that move. Distinct moves have distinct labels. Thus the set of moves injects into $I\times\N$, whose cardinality is at most $\kappa$. A run of length at least $\kappa^+$ would have more than $\kappa$ moves, an impossibility.

Each $G_i$ has a set of reachable forms, because taking its transitive closure along its set-valued options is a set construction. Their product is a set of positions. Well-order the set of legal moves from all positions, using Choice. Starting after a given history, at every successor choose the first legal move for the designated player, if any; at limits use the first lemma. If this recursion never stopped below $\kappa^+$ it would construct an impossible run of that length. It therefore stops at a maximal alternating history. All histories have ordinal domains below $\kappa^+$ and values in a fixed set; their union over these domains is a set. Strategies, being appropriate functions on subsets of that set, are also sets. If $I$ is countable, $\kappa^+=\omega_1$, proving the claimed countability.
\end{proof}

\begin{proposition}[No countable bound even for one fixed sum]
Let every coordinate in a countably infinite index set be the one-move impartial game $*=\{0\mid0\}$. For every countably infinite ordinal $\alpha$ there is a maximal alternating play of length exactly $\alpha$ on this same initial $\Delta$-sum. In particular no countable ordinal bounds all its play lengths.
\end{proposition}
\begin{proof}
Choose a bijection $f:\alpha\to\N$. At move time $\beta<\alpha$, remove the star in coordinate $f(\beta)$. It is available because it was not used earlier, and either player may remove it. At limit stages, each removed coordinate stays zero and every other one stays a star, as required. For each $\beta<\alpha$ some stars remain, because the bijection has not yet used $f(\beta)$. Thus play cannot stop earlier. At time $\alpha$ every coordinate is zero, and the player due to move loses. This is a maximal play with the prescribed length.
\end{proof}

\section*{A rigorously determined infinite class}
The conjugate $-H$ is obtained from a game form $H$ by interchanging Left and Right options recursively. The following hypothesis is about an exact form isomorphism, including the options, rather than Conway equivalence.

\begin{theorem}[Paired conjugate sums are second-player wins]
Suppose the nonzero coordinates of a $\Delta$-sum can be partitioned into disjoint pairs $\{i,j\}$ such that the two initial forms in each pair are conjugate. Coordinates equal to the terminal form $0$ may be left unpaired; here zero refers to that form, not merely Conway equivalence to zero. Then for either designation of First, Second has a winning strategy under the stated transfinite limit rule.
\end{theorem}
\begin{proof}
Fix the pairing and the conjugating form isomorphisms. Maintain the invariant that immediately before every First move the two forms in each pair are conjugate. If First changes $H$ in one coordinate by a move of First's color, the conjugating correspondence supplies a legal move of Second's color in the other coordinate. Second plays that move. The two resulting forms are again conjugate, and all other pairs were unchanged.

It remains to verify the invariant at limits, which is the part not supplied by the usual finite argument. Fix a limit stage $\lambda$ and one pair. By finite coordinate activity, only finitely many moves before $\lambda$ have affected either coordinate in that pair. Every First move before $\lambda$ has its immediately following response also before $\lambda$, since $\lambda$ is a limit ordinal. Thus, after the last such completed response affecting the pair, its forms are conjugate and remain unchanged up to $\lambda$. Their limit forms are therefore conjugate. This holds for every pair, so the invariant is restored at the limit, when First is again due to move.

Second can never be the first player without a move: after each First move the conjugate response exists. A maximal play compatible with this strategy exists: perform the same transfinite extension recursion as in the cardinal-bound proposition, using the prescribed response on Second turns and any well-ordered legal choice on First turns. The cardinal bound forces this compatible recursion to terminate. At termination First, rather than Second, must be unable to move. Hence Second wins every maximal compatible play, as required by the definition of a winning strategy.
\end{proof}

\begin{corollary}[Absorption of one extra star in this example]
A countably infinite sum of stars is a second-player win. Adding one more star gives another second-player win, although one star by itself is a first-player win.
\end{corollary}
\begin{proof}
The star is its own conjugate, and any countably infinite set can be partitioned into pairs. Adding one coordinate leaves the index set countably infinite, so the theorem applies in both cases. A single star has one legal move to zero, which wins for First.
\end{proof}

\section*{What fails in a tempting determinacy proof}
The cardinal bound says that each countable-sum play ends before $\omega_1$; it does not provide a countable ordinal decreasing at every move. The fixed all-star example has arbitrarily long finite prefixes and has plays of every countably infinite length. In particular its ordinary finite-history extension tree has an infinite branch, so it is not a well-founded tree in the sense needed for ordinary backward recursion.

There is no contradiction with the paired-sum theorem: that theorem verifies a strategy at every limit directly, without a countdown rank. Nor does exhaustion of each actual play prove that one of the two players has a strategy controlling all possible plays. General unpaired sums, including the countable case, remain unresolved here. No determinacy axiom stronger than ZFC is invoked or claimed necessary.

\begin{thebibliography}{9}
\bibitem{lipparini} P. Lipparini, \emph{Infinite sums of combinatorial games (Dadaist games)}, arXiv:2505.00614v1 (2025), Section 3, especially Definitions 3.1 and 3.3, Problem 3.4, and Remark 3.2. \url{https://arxiv.org/html/2505.00614v1}. The primary source and the convention that First restarts at limits were checked on 9 October 2026. The arguments above use only these rules and the explicitly stated ZFC set constructions.
\end{thebibliography}

\end{document}
